\documentclass[12pt]{article}
\usepackage[utf8]{inputenc}
\usepackage[T1]{fontenc}
\usepackage{hyperref}
\title{Kapin Language Resources for K-12 Curriculum Development}
\author{Compiled by Claude Sonnet 5 (claude-sonnet-5) for Andr\'as Kornai}
\date{}
\begin{document}
\maketitle
\section{Introduction}
Kapin (ISO 639-3 tbx) is an Oceanic (Austronesian) language of Morobe Province, Papua New Guinea, with about 2,400 speakers (1979 census per Wikipedia). Joshua Project lists it under the Middle Watut / Sambio people group (alternate names Sambio, Tayek), with a much larger population figure; the sources differ in method and year. Little material exists online.
\section{Primary Educational Resources}
No primers, textbooks or school programmes were found.
\section{Linguistic and Reference Materials}
Wikipedia, Kapin language. \url{https://en.wikipedia.org/wiki/Kapin_language}. No grammar or dictionary was found.
\section{Religious and Community Resources}
\begin{itemize}
\item Bible portions in Kapin on ebible.org. \url{https://ebible.org/details.php?id=tbx}
\item The same on png.bible. \url{https://png.bible/details.php?id=tbx}
\item Joshua Project, Middle Watut, Sambio people group. \url{https://joshuaproject.net/people_groups/13710}
\end{itemize}
\section{Digital and Audio Resources}
Global Recordings Network page for Kapin (audio/gospel recordings; content not inspected). \url{https://globalrecordings.net/en/language/tbx}
\section{Curriculum Development Context}
The available materials are mainly Bible portions and recordings, which give some written text and audio in the language and indicate an orthography exists. No secular educational or linguistic descriptions were found, so a curriculum would need new community-led materials.
\end{document}